\section{问题分析}

% 写作提示：逐问说明输入、输出、约束、数据特点、关键难点、题型判断、拟采用方法及适用理由，
% 并交代各问之间的依赖关系。本节重点回答”为什么这样做”，通常不展开公式或提前给出结果。

\subsection{问题一的分析}

问题一的核心挑战是\textbf{多模态时序异构性}：文本模态以词为单位（离散、稀疏），语音模态以固定帧率采样（连续、密集），视觉模态以视频帧为单位（中等密度），三者在时间粒度、采样频率和序列长度上均不一致。直接拼接会导致维度爆炸和时序错位。

\textbf{数据特点}：附件2提供的aligned\_50.pkl包含4850个样本，每个样本的文本特征维度768（BERT编码），视觉特征维度35（面部动作单元），音频特征维度74（COVAREP声学特征）。未对齐数据（unaligned\_50.pkl）中每个模态是变长时序，需要统一为固定长度表示。

\textbf{拟采用方法}：采用统计特征聚合策略，从时序维度提取均值、标准差、最大值、最小值、中位数、四分位距等统计量，将变长序列映射为固定维度特征向量。该方法计算高效、鲁棒性强，能保留时序分布信息而不依赖复杂的序列模型。

\subsection{问题二的分析}

问题二的核心难点是\textbf{模态局部缺失的不确定性}：训练数据中部分时刻的某些模态为全零（缺失标记），测试时缺失模式未知。传统固定权重融合难以适应动态缺失场景，需要模型具备模态感知和自适应能力。

\textbf{数据特点}：附件3提供30个缺失样本，缺失率从30\%到70\%不等，缺失模式包括单模态完全缺失、多模态交替缺失、随机时刻缺失等。情感标签为连续回归值[-3, 3]，需要保持预测稳定性。

\textbf{拟采用方法}：采用梯度提升树（Gradient Boosting）算法，训练时引入随机模态缺失机制（Dropout）增强鲁棒性。梯度提升树能够自动学习特征重要性和交互模式，对缺失值不敏感，且可通过特征重要性分析量化各模态贡献度。该方法避免了深度学习的过拟合风险，适合中等规模数据集。

\subsection{问题三的分析}

问题三的核心需求是\textbf{预测结果的可追溯性}：不仅要给出情感分数，还需要定位关键时刻、识别重要模态、回溯到原始证据。这要求模型具备细粒度的特征贡献度量化能力。

\textbf{数据特点}：附件4提供20个专项样本，需要针对每个预测结果生成可解释性报告，包括模态权重分布、时序关键帧定位、特征贡献度排序等。

\textbf{拟采用方法}：基于SHAP（SHapley Additive exPlanations）值方法，量化每个特征对预测结果的边际贡献。SHAP值源自博弈论的Shapley值，具有一致性、局部准确性和可加性等理论保证。通过SHAP值可视化，可以生成模态贡献度柱状图、时序注意力热图、关键特征排序等多维度解释，支持决策审查和用户信任。

\subsection{三问之间的依赖关系}

问题一为后续建模提供统一的特征表示，是问题二和问题三的数据基础。问题二的鲁棒预测模型可直接用于问题三的可解释性分析，通过SHAP值方法在问题二模型基础上进行后处理即可实现。三个问题构成”特征工程→鲁棒建模→可解释分析”的完整技术链条。

